\section*{Sammendrag}

I dette prosjektet ble det utviklet og analysert en RF-krets bestående av en NE555-timer, en transistor, en LC-resonanskrets og en antenne. Målet var å generere et høyfrekvent signal i nærheten av FM-båndet og undersøke hvordan teoretiske beregninger samsvarer med praktiske målinger.
\\[1em]
Det ble gjennomført teoretiske analyser av kretsens ulike deler, inkludert beregninger av timerfrekvens, transistorens arbeidspunkt og resonansfrekvensen til LC-kretsen. Beregningene viste at kretsen i teorien skulle kunne generere signaler innenfor FM-båndet. Kretsen ble først testet på breadboard, hvor den ga merkbar forstyrrelse på en FM-radio.
\\[1em]
Etter overføring til PCB fungerte kretsen imidlertid ikke som forventet. Feilsøking ved hjelp av spennings- og frekvensmålinger viste store avvik mellom teori og praksis. Særlig ble det observert svært lav kollektorspenning på transistoren og frekvenser langt under det forventede området. Målingene indikerte at dårlig elektrisk kontakt i loddepunkter og tilkoblinger til spoler og bryter var de mest sannsynlige årsakene til feilen.
\\[1em]
Rapporten viser at selv om de teoretiske beregningene var konsistente med ønsket funksjon, kan praktiske forhold som loddekvalitet, komponentmontering og fysiske koblinger være avgjørende for om en RF-krets fungerer som tiltenkt.
